\section{Conclusioni e sviluppi futuri}

Le quattro domande di ricerca ricevono una risposta convergente. Rispetto a RQ1, gli account successivamente sottoposti a takedown ricevono piu' blocchi dei controlli e il segnale e' fortemente concentrato negli ultimi giorni, soprattutto nel giorno immediatamente precedente all'evento. Rispetto a RQ2, il potere informativo dei blocchi dipende dall'esposizione: funziona meglio per account con post e/o follow ricevuti, mentre e' debole per il corner 00. Rispetto a RQ3, una Random Forest basata sui blocker unici consente una discriminazione parziale ma non completa, con F1 out-of-fold pari a \num{0.6878} e ROC AUC pari a \num{0.7604}. Rispetto a RQ4, labels ufficiali post-level, labeler di terze parti e modlist hanno coverage troppo bassa per sostenere un uso predittivo sistematico, anche quando il lift e' elevato.

Il messaggio centrale e' che i blocchi ricevuti costituiscono il segnale comunitario piu' informativo tra quelli analizzati, ma il loro potere predittivo dipende dall'esposizione pubblica dell'account. Il segnale e' forte e temporalmente concentrato per gli account che hanno avuto il tempo di interagire con la comunita', mentre non puo' spiegare i takedown rapidissimi degli account privi di post e follow ricevuti. Labels e modlist presentano talvolta lift elevati, ma una coverage insufficiente per sostenere una predizione sistematica.

Il contributo principale e' distinguere empiricamente due famiglie di takedown: quelli che lasciano segnali comunitari pubblici nella finestra pre-evento e quelli che avvengono prima che tali segnali possano formarsi. Questa distinzione aiuta a interpretare la moderazione su Bluesky come interazione tra segnali comunitari osservabili e sistemi ufficiali account-level non completamente osservabili.

Gli sviluppi futuri includono l'estensione a piu' mesi, una validazione temporale su periodi successivi, grouped cross-validation per DID, validazione indipendente del modello, curve precision-recall, calibrazione delle probabilita' e sensitivity analysis su finestre di 3, 5, 7, 10 e 14 giorni. Ulteriori analisi potrebbero modellare il tempo al takedown con approcci di sopravvivenza, studiare il grafo dei blocker, includere centralita' e reputazione dei blocker, analizzare separatamente il corner 00 e classificare le cause dei takedown. Per l'interpretabilita' del modello, permutation importance e, in futuro, SHAP potrebbero essere applicati come tecniche aggiuntive; SHAP non e' stato usato nell'analisi corrente.
